% ===========================================================================
% Section I -- Introduction
% Owner: Member 1
% ===========================================================================
\IEEEPARstart{R}{olling} element bearings support rotating shafts in motors,
pumps, turbines, gearboxes and many other machines. Their condition matters
for machine availability, because localized surface damage can develop into
severe mechanical failure when it is not found early. A defect on an inner
race, outer race or rolling element produces a short-duration contact
disturbance each time the rolling contact passes the damaged region. These
impacts excite resonances in the bearing and in the supporting structure, and
an accelerometer mounted on or near the housing observes them as vibration.
The diagnostic information sits in the \emph{repetition rate} of the impacts,
and it is recovered by demodulating the resonance that carries them, which is
the principle behind envelope analysis~\cite{randall2011,kim2020}.

The measured vibration is not stationary across operating conditions. Load
changes the contact force and the strength of the defect-induced impacts, and
it also changes the overall vibration and noise floor. Effective shaft speed
differs between load conditions, and local slip makes the actual impact
repetition rate depart slightly from the ideal value predicted by bearing
geometry. A fixed threshold on raw amplitude or on the spectrum is therefore
not expected to behave consistently at all loads.

The approach taken here is a pure DSP framework, not a model trained on
labelled examples. Every stage can be inspected on its own: the
characteristic frequencies, the filters, the transforms and the detection
thresholds are all explicit. Interpretability of this kind does not guarantee
generalization. Bearing geometry, shaft speed, sensor placement and
validation data are still required.

This paper applies two classical operations where they matter most. The
filter band is chosen by kurtosis, which measures impulsiveness and therefore
locates the resonance excited by the defect~\cite{antoni2007}; the decision
threshold is estimated \emph{locally} in the envelope spectrum with an order
statistic, which is insensitive to the impulsive outliers that corrupt a
moving average. The resulting detector contains no learned parameters and can
be implemented on a fixed-point microcontroller.

The contributions of this paper are:
\begin{itemize}
  \item a Kurtosis-guided resonance band selection procedure embedded in a
        cascade-SOS zero-phase Butterworth filter, compared against a
        linear-phase Kaiser FIR design on the basis of pole--zero location,
        magnitude response and group delay;
  \item a Hilbert squared-envelope demodulator with calibrated Hann-windowed
        FFT, in which squaring is shown by the modulation theorem to shift the
        resonance to baseband without the intermodulation produced by a
        square-root envelope;
  \item a Median--MAD CFAR threshold that tracks the load-dependent noise
        floor, with a slip-tolerant harmonic comb, evaluated against static
        thresholding and cell-averaging CFAR on multi-load data.
\end{itemize}
